\documentclass[UTF8,zihao=-4,fontset=none]{ctexart}

\usepackage[a4paper,top=2.45cm,bottom=2.2cm,left=2.55cm,right=2.55cm,headheight=15pt]{geometry}
\usepackage{fontspec}
\usepackage{xeCJK}
\usepackage{unicode-math}
\usepackage{amsmath}
\usepackage{graphicx}
\usepackage{booktabs}
\usepackage{array}
\usepackage{tabularx}
\usepackage{multirow}
\usepackage{float}
\usepackage{caption}
\usepackage{fancyhdr}
\usepackage{enumitem}
\usepackage[hidelinks,unicode]{hyperref}
\usepackage{bookmark}

\setmainfont{Times New Roman}
\setsansfont{Times New Roman}
\setmonofont{Consolas}
\setCJKmainfont{SimSun}[AutoFakeBold=true,AutoFakeSlant=true]
\setCJKsansfont{SimHei}[AutoFakeBold=true]
\setmathfont{TeX Gyre Termes Math}

\ctexset{
  section = {
    format = \zihao{3}\bfseries,
    beforeskip = 1.0em,
    afterskip = 0.8em
  },
  subsection = {
    format = \zihao{4}\bfseries,
    beforeskip = 0.8em,
    afterskip = 0.5em
  }
}

\renewcommand{\contentsname}{目\quad 录}
\renewcommand{\arraystretch}{1.18}
\renewcommand{\thetable}{\thesection-\arabic{table}}
\makeatletter
\@addtoreset{table}{section}
\makeatother
\setlength{\parindent}{2em}
\setlength{\parskip}{0pt}
\linespread{1.18}
\captionsetup{font=small,labelsep=quad}
\newcommand{\stepheading}[1]{\vspace{0.6em}\noindent\textbf{#1}\par}

\newcommand{\studentname}{张天羽}
\newcommand{\studentid}{3240103613}
\newcommand{\coursename}{系统建模与仿真}
\newcommand{\quiztitle}{Quiz 3}

\fancypagestyle{coverstyle}{
  \fancyhf{}
  \fancyhead[L]{\zihao{5}\studentname\ \studentid}
  \fancyhead[C]{\zihao{5}\coursename}
  \fancyhead[R]{\zihao{5}\quiztitle}
  \renewcommand{\headrulewidth}{0.4pt}
}

\fancypagestyle{mainstyle}{
  \fancyhf{}
  \fancyhead[L]{\zihao{5}\studentname\ \studentid}
  \fancyhead[C]{\zihao{5}\coursename}
  \fancyhead[R]{\zihao{5}\nouppercase{\leftmark}}
  \fancyfoot[C]{\thepage}
  \renewcommand{\headrulewidth}{0.4pt}
}

\renewcommand{\sectionmark}[1]{\markboth{\thesection\ #1}{}}

\begin{document}
\zihao{-4}

\pagenumbering{gobble}
\thispagestyle{coverstyle}
\begin{center}
  \vspace*{1.1cm}
  {\zihao{2}\bfseries \quiztitle\par}
  \vspace{0.7cm}
  {\zihao{4}\studentname\par}
\end{center}
\vspace{0.4cm}
\tableofcontents

\clearpage
\pagestyle{mainstyle}
\pagenumbering{arabic}

\section{第一题}

\subsection{问题叙述}

1. A farmer has 30 acres on which to grow tomatoes and corn. Each 100 bushels of tomatoes require 1000 gallons of water and 5 acres of land. Each 100 bushels of corn require 6000 gallons of water and 2.5 acres of land. Labor costs are \$1 per bushel for both corn and tomatoes. The farmer has available 30,000 gallons of water and \$750 in capital. He knows that he cannot sell more than 500 bushels of tomatoes or 475 bushels of corn. He estimates a profit of \$2 on each bushel of tomatoes and \$3 on each bushel of corn.

\begin{enumerate}[label=\alph*.,leftmargin=2.3em]
  \item How many bushels of each should he raise to maximize profits?
  \item Next, assume that the farmer has the opportunity to sign a nice contract with a grocery store to grow and deliver at least 300 bushels of tomatoes and at least 500 bushels of corn. Should the farmer sign the contract? Support your recommendation.
  \item Now assume that the farmer can obtain an additional 10,000 gallons of water for a total cost of \$50. Should he obtain the additional water? Support your recommendation.
\end{enumerate}

\subsection{模型建立}

令
\[
    x_1=\text{tomatoes 的产量，单位为 }100\text{ bushels},
    \qquad
    x_2=\text{corn 的产量，单位为 }100\text{ bushels}.
\]
这样设置变量后，水、土地、资本以及销售上限都可以直接由题目数据写出。

每 \(100\) bushels tomatoes 可获得利润 \(200\) dollars，每 \(100\) bushels corn 可获得利润 \(300\) dollars。以 \(100\) dollars 为利润单位，目标函数为
\[
    \max z=2x_1+3x_2. \tag{1-1}
\]
水资源约束为
\[
    1000x_1+6000x_2\le 30000,
\]
即
\[
    x_1+6x_2\le 30. \tag{1-2}
\]
土地约束为
\[
    5x_1+2.5x_2\le 30,
\]
即
\[
    2x_1+x_2\le 12. \tag{1-3}
\]
资本约束为
\[
    100x_1+100x_2\le 750,
\]
即
\[
    x_1+x_2\le 7.5. \tag{1-4}
\]
销售上限约束为
\[
    x_1\le 5,\qquad x_2\le 4.75. \tag{1-5}
\]
因此，第一题(a)的线性规划模型为
\[
    \begin{aligned}
    \max\quad & z=2x_1+3x_2,\\
    \mathrm{s.t.}\quad
      &x_1+6x_2\le 30,\\
      &2x_1+x_2\le 12,\\
      &x_1+x_2\le 7.5,\\
      &x_1\le 5,\\
      &x_2\le 4.75,\\
      &x_1,x_2\ge 0.
    \end{aligned} \tag{1-6}
\]

\subsection{单纯形法推导}

单纯形法的基本思想是：线性规划的最优解若存在，就可以在可行域的某个顶点上取得。因此不必在整个可行域中盲目寻找，而是从一个容易得到的顶点出发，沿着能够使目标函数变好的边移动到相邻顶点。每移动一次，就把一个原来取 \(0\) 的非基变量放入基中，同时让一个原来的基变量降到 \(0\) 并离开基。这样既保持约束方程成立，又始终不破坏非负性。

下面先不使用单纯形表，而是直接从基变量与非基变量之间的关系出发进行换基迭代。每一步都需要回答两个问题：第一，哪个非基变量增加后能使目标函数增加，它就应当入基；第二，这个变量最多能增加到多少，最先被压到 \(0\) 的基变量就必须出基。

\stepheading{化为标准形式}

引入松弛变量 \(x_3,x_4,x_5,x_6,x_7\)，得到
\[
    \begin{aligned}
    \max\quad & z=2x_1+3x_2,\\
    \mathrm{s.t.}\quad
      &x_1+6x_2+x_3=30,\\
      &2x_1+x_2+x_4=12,\\
      &x_1+x_2+x_5=7.5,\\
      &x_1+x_6=5,\\
      &x_2+x_7=4.75,\\
      &x_1,x_2,x_3,x_4,x_5,x_6,x_7\ge 0.
    \end{aligned} \tag{1-7}
\]
其系数矩阵为
\[
    A=(P_1,P_2,P_3,P_4,P_5,P_6,P_7)
    =
    \begin{bmatrix}
      1&6&1&0&0&0&0\\
      2&1&0&1&0&0&0\\
      1&1&0&0&1&0&0\\
      1&0&0&0&0&1&0\\
      0&1&0&0&0&0&1
    \end{bmatrix}. \tag{1-8}
\]
显然，列向量组
\[
    B=(P_3,P_4,P_5,P_6,P_7)=I
\]
构成单位矩阵的基。取 \(x_3,x_4,x_5,x_6,x_7\) 为初始基变量，取 \(x_1,x_2\) 为非基变量。将所有基变量用非基变量表示，可得
\[
    \begin{cases}
    x_3=30-x_1-6x_2,\\
    x_4=12-2x_1-x_2,\\
    x_5=7.5-x_1-x_2,\\
    x_6=5-x_1,\\
    x_7=4.75-x_2.
    \end{cases} \tag{1-9}
\]
令非基变量 \(x_1=x_2=0\)，得初始基可行解
\[
    X^{(0)}=(0,0,30,12,7.5,5,4.75)^{\mathrm T},\qquad z^{(0)}=0. \tag{1-10}
\]
这个解的含义是：暂时不生产 tomatoes 和 corn，所有资源约束的剩余量都由松弛变量表示。虽然这个方案显然不是利润最大的方案，但它满足所有约束，而且由单位矩阵基直接得到，所以可以作为单纯形法的起点。

由
\[
    z=2x_1+3x_2 \tag{1-11}
\]
可见，在当前基可行解附近，若把 \(x_1\) 从 \(0\) 增加一个很小的正数 \(\Delta\)，目标函数增加 \(2\Delta\)；若把 \(x_2\) 从 \(0\) 增加同样的 \(\Delta\)，目标函数增加 \(3\Delta\)。两者都能提高利润，而 \(x_2\) 的单位增益更大，所以按“最大正系数”规则选择
\[
    x_2
\]
为入基变量。这里需要区分：目标函数中较大的正系数决定的是入基变量，而出基变量不能由目标函数系数决定，必须由可行性条件决定。

\stepheading{第一次换基迭代}

保持 \(x_1=0\)，让 \(x_2\) 从 \(0\) 开始增加。由式(1-9)可得
\[
    \begin{cases}
    x_3=30-6x_2\ge 0,\\
    x_4=12-x_2\ge 0,\\
    x_5=7.5-x_2\ge 0,\\
    x_6=5\ge 0,\\
    x_7=4.75-x_2\ge 0.
    \end{cases} \tag{1-12}
\]
为了保持可行解，所有变量都必须非负。式(1-12)中的每一行都给出了 \(x_2\) 可以增大的上界：
\[
    x_3\ge0\Rightarrow x_2\le \frac{30}{6},\quad
    x_4\ge0\Rightarrow x_2\le \frac{12}{1},\quad
    x_5\ge0\Rightarrow x_2\le \frac{7.5}{1},\quad
    x_7\ge0\Rightarrow x_2\le \frac{4.75}{1}.
\]
\(x_6=5\) 与 \(x_2\) 无关，所以不限制 \(x_2\) 的增加。若 \(x_2\) 超过这些上界中的任意一个，相应的基变量就会变成负数，解将不再可行；若停在最小上界处，则正好有一个基变量变为 \(0\)，可以得到相邻的顶点。因此步长取
\[
    \theta=\min\left\{\frac{30}{6},\frac{12}{1},\frac{7.5}{1},\frac{4.75}{1}\right\}=4.75. \tag{1-13}
\]
当 \(x_2=4.75\) 时，\(x_7=0\)，所以 \(x_7\) 为出基变量。

由
\[
    x_2+x_7=4.75
\]
得
\[
    x_2=4.75-x_7. \tag{1-14}
\]
这一步的意义是把 \(x_2\) 改写为新的基变量，把 \(x_7\) 改写为新的非基变量。之后令新的非基变量为 \(0\) 时，自动得到换基后的相邻顶点。

代入式(1-9)，得
\[
    \begin{cases}
    x_3=1.5-x_1+6x_7,\\
    x_4=7.25-2x_1+x_7,\\
    x_5=2.75-x_1+x_7,\\
    x_6=5-x_1,\\
    x_2=4.75-x_7.
    \end{cases} \tag{1-15}
\]
目标函数变为
\[
    z=2x_1+3(4.75-x_7)=14.25+2x_1-3x_7. \tag{1-16}
\]
令新的非基变量 \(x_1=x_7=0\)，得到
\[
    X^{(1)}=(0,4.75,1.5,7.25,2.75,5,0)^{\mathrm T},\qquad z^{(1)}=14.25. \tag{1-17}
\]
式(1-16)说明，在新的顶点处，若增加 \(x_1\)，目标函数按 \(2\) 的速率增加；若增加 \(x_7\)，目标函数按 \(-3\) 的速率变化，即会降低利润。因此只有 \(x_1\) 是能继续改善目标函数的方向，选择 \(x_1\) 为新的入基变量。

\stepheading{第二次换基迭代}

保持 \(x_7=0\)，让 \(x_1\) 从 \(0\) 开始增加。由式(1-15)可得
\[
    \begin{cases}
    x_3=1.5-x_1\ge 0,\\
    x_4=7.25-2x_1\ge 0,\\
    x_5=2.75-x_1\ge 0,\\
    x_6=5-x_1\ge 0,\\
    x_2=4.75\ge 0.
    \end{cases} \tag{1-18}
\]
同样地，入基变量 \(x_1\) 的增加不能使任何基变量变成负数。由式(1-18)可知
\[
    x_3\ge0\Rightarrow x_1\le1.5,\quad
    x_4\ge0\Rightarrow x_1\le\frac{7.25}{2},\quad
    x_5\ge0\Rightarrow x_1\le2.75,\quad
    x_6\ge0\Rightarrow x_1\le5.
\]
\(x_2=4.75\) 不随 \(x_1\) 变化，不会限制此次移动。为了在保持可行的前提下沿改善方向走到下一个顶点，应取这些上界中的最小值：
\[
    \theta=\min\left\{1.5,\frac{7.25}{2},2.75,5\right\}=1.5. \tag{1-19}
\]
当 \(x_1=1.5\) 时，\(x_3=0\)，所以 \(x_3\) 为出基变量。

由
\[
    x_3=1.5-x_1+6x_7
\]
解出
\[
    x_1=1.5+6x_7-x_3. \tag{1-20}
\]
于是 \(x_1\) 进入基，\(x_3\) 离开基。把这个关系代回其余方程，就能把所有新的基变量继续表示成新的非基变量。

代入式(1-15)，得
\[
    \begin{cases}
    x_1=1.5+6x_7-x_3,\\
    x_4=4.25+2x_3-11x_7,\\
    x_5=1.25+x_3-5x_7,\\
    x_6=3.5+x_3-6x_7,\\
    x_2=4.75-x_7.
    \end{cases} \tag{1-21}
\]
目标函数变为
\[
    z=14.25+2(1.5+6x_7-x_3)-3x_7=17.25-2x_3+9x_7. \tag{1-22}
\]
令新的非基变量 \(x_3=x_7=0\)，得到
\[
    X^{(2)}=(1.5,4.75,0,4.25,1.25,3.5,0)^{\mathrm T},\qquad z^{(2)}=17.25. \tag{1-23}
\]
式(1-22)中，\(x_3\) 的系数为 \(-2\)，增加 \(x_3\) 会降低目标函数；\(x_7\) 的系数为 \(9\)，增加 \(x_7\) 会提高目标函数。因此选择 \(x_7\) 为新的入基变量。虽然 \(x_7\) 本身是 corn 销售上限约束的松弛变量，但在当前基表示下，增加 \(x_7\) 会同时调整 \(x_1,x_2\) 等变量，整体效果是使目标函数增加。

\stepheading{第三次换基迭代}

保持 \(x_3=0\)，让 \(x_7\) 从 \(0\) 开始增加。由式(1-21)可得
\[
    \begin{cases}
    x_1=1.5+6x_7\ge 0,\\
    x_4=4.25-11x_7\ge 0,\\
    x_5=1.25-5x_7\ge 0,\\
    x_6=3.5-6x_7\ge 0,\\
    x_2=4.75-x_7\ge 0.
    \end{cases} \tag{1-24}
\]
式(1-24)继续体现可行性对步长的限制。由于
\[
    x_1=1.5+6x_7\ge0
\]
在 \(x_7\ge0\) 时总成立，所以 \(x_1\) 不限制 \(x_7\) 增大。其余基变量给出
\[
    x_4\ge0\Rightarrow x_7\le\frac{4.25}{11},\quad
    x_5\ge0\Rightarrow x_7\le\frac{1.25}{5},\quad
    x_6\ge0\Rightarrow x_7\le\frac{3.5}{6},\quad
    x_2\ge0\Rightarrow x_7\le4.75.
\]
仍然取最小上界，才能在不破坏非负性的条件下到达相邻顶点：
\[
    \theta=\min\left\{\frac{4.25}{11},\frac{1.25}{5},\frac{3.5}{6},4.75\right\}=0.25. \tag{1-25}
\]
当 \(x_7=0.25\) 时，\(x_5=0\)，所以 \(x_5\) 为出基变量。

由
\[
    x_5=1.25+x_3-5x_7
\]
解出
\[
    x_7=0.25+\frac{1}{5}x_3-\frac{1}{5}x_5. \tag{1-26}
\]
这表示 \(x_7\) 成为新的基变量，\(x_5\) 成为新的非基变量。继续代入后，可以得到新的顶点表达式。

代入式(1-21)，得
\[
    \begin{cases}
    x_7=0.25+\dfrac{1}{5}x_3-\dfrac{1}{5}x_5,\\[0.4em]
    x_1=3+\dfrac{1}{5}x_3-\dfrac{6}{5}x_5,\\[0.4em]
    x_4=1.5-\dfrac{1}{5}x_3+\dfrac{11}{5}x_5,\\[0.4em]
    x_6=2-\dfrac{1}{5}x_3+\dfrac{6}{5}x_5,\\[0.4em]
    x_2=4.5-\dfrac{1}{5}x_3+\dfrac{1}{5}x_5.
    \end{cases} \tag{1-27}
\]
目标函数变为
\[
    z=17.25-2x_3+9\left(0.25+\frac{1}{5}x_3-\frac{1}{5}x_5\right)
     =19.5-\frac{1}{5}x_3-\frac{9}{5}x_5. \tag{1-28}
\]
令新的非基变量 \(x_3=x_5=0\)，得到
\[
    X^{(3)}=(3,4.5,0,1.5,0,2,0.25)^{\mathrm T},\qquad z^{(3)}=19.5. \tag{1-29}
\]
式(1-28)中，当前非基变量为 \(x_3,x_5\)。由于它们都受到非负约束，只能从 \(0\) 向正方向增加；而二者在目标函数中的系数分别为 \(-1/5\) 和 \(-9/5\)，增加任何一个都会使 \(z\) 不增反降。因此，从当前顶点出发已经没有任何相邻方向可以继续提高目标函数，\(X^{(3)}\) 即为最优解。于是
\[
    x_1=3,\qquad x_2=4.5.
\]
由于 \(x_1,x_2\) 的单位为 \(100\) bushels，第一题(a)的答案为
\[
    \boxed{300\text{ bushels tomatoes},\qquad 450\text{ bushels corn}.}
\]
最大利润为
\[
    \boxed{1950\text{ dollars}.}
\]

\subsection{单纯形表}

下面用单纯形表重复同一计算过程。表中最后一行为检验数
\[
    \lambda_j=c_j-\sum_i C_{B_i}a_{ij}.
\]
检验数的作用与前面目标函数中非基变量的系数相同：正检验数表示该列变量入基后可以提高目标函数；若有多个正检验数，取最大者作为入基变量，可以使本次单位步长的改进最大。确定入基列后，再只在该列正系数所在的行计算 \(b_i/a_{ij}\)。这些比值就是入基变量继续增加时各基变量能保持非负的最大允许值，最小正比值所在行首先达到 \(0\)，因此该行基变量出基。

\begin{table}[H]
\centering
\zihao{-4}
\setlength{\tabcolsep}{2.2pt}
\caption{初始单纯形表}
\begin{tabular}{ccccccccccc}
\toprule
\multicolumn{3}{c}{$c_j\rightarrow$} & \(2\) & \(3\) & \(0\) & \(0\) & \(0\) & \(0\) & \(0\) & \multirow{2}{*}{\(\theta_i\)}\\
\(C_B\) & 基 & \(b\) & \(x_1\) & \(x_2\) & \(x_3\) & \(x_4\) & \(x_5\) & \(x_6\) & \(x_7\) &\\
\midrule
\(0\) & \(x_3\) & \(30\)   & \(1\) & \(6\) & \(1\) & \(0\) & \(0\) & \(0\) & \(0\) & \(30/6=5\)\\
\(0\) & \(x_4\) & \(12\)   & \(2\) & \(1\) & \(0\) & \(1\) & \(0\) & \(0\) & \(0\) & \(12/1=12\)\\
\(0\) & \(x_5\) & \(7.5\)  & \(1\) & \(1\) & \(0\) & \(0\) & \(1\) & \(0\) & \(0\) & \(7.5/1=7.5\)\\
\(0\) & \(x_6\) & \(5\)    & \(1\) & \(0\) & \(0\) & \(0\) & \(0\) & \(1\) & \(0\) & --\\
\(0\) & \(x_7\) & \(4.75\) & \(0\) & \([1]\) & \(0\) & \(0\) & \(0\) & \(0\) & \(1\) & \(4.75/1=4.75\)\\
\multicolumn{3}{c}{检验数} & \(2\) & \(3\) & \(0\) & \(0\) & \(0\) & \(0\) & \(0\) &\\
\bottomrule
\end{tabular}
\end{table}

因为
\[
    \lambda_2=3
\]
为最大正检验数，所以 \(x_2\) 入基。按 \(\theta\) 规则，\(x_7\) 出基。换基运算为
\[
    R_{x_2}=R_{x_7},\qquad
    R_{x_3}\leftarrow R_{x_3}-6R_{x_2},\qquad
    R_{x_4}\leftarrow R_{x_4}-R_{x_2},\qquad
    R_{x_5}\leftarrow R_{x_5}-R_{x_2}.
\]

\begin{table}[H]
\centering
\zihao{-4}
\setlength{\tabcolsep}{2.2pt}
\caption{第一次换基后的单纯形表}
\begin{tabular}{ccccccccccc}
\toprule
\multicolumn{3}{c}{$c_j\rightarrow$} & \(2\) & \(3\) & \(0\) & \(0\) & \(0\) & \(0\) & \(0\) & \multirow{2}{*}{\(\theta_i\)}\\
\(C_B\) & 基 & \(b\) & \(x_1\) & \(x_2\) & \(x_3\) & \(x_4\) & \(x_5\) & \(x_6\) & \(x_7\) &\\
\midrule
\(0\) & \(x_3\) & \(1.5\)  & \([1]\) & \(0\) & \(1\) & \(0\) & \(0\) & \(0\) & \(-6\) & \(1.5/1=1.5\)\\
\(0\) & \(x_4\) & \(7.25\) & \(2\) & \(0\) & \(0\) & \(1\) & \(0\) & \(0\) & \(-1\) & \(7.25/2=3.625\)\\
\(0\) & \(x_5\) & \(2.75\) & \(1\) & \(0\) & \(0\) & \(0\) & \(1\) & \(0\) & \(-1\) & \(2.75/1=2.75\)\\
\(0\) & \(x_6\) & \(5\)    & \(1\) & \(0\) & \(0\) & \(0\) & \(0\) & \(1\) & \(0\) & \(5/1=5\)\\
\(3\) & \(x_2\) & \(4.75\) & \(0\) & \(1\) & \(0\) & \(0\) & \(0\) & \(0\) & \(1\) & --\\
\multicolumn{3}{c}{检验数} & \(2\) & \(0\) & \(0\) & \(0\) & \(0\) & \(0\) & \(-3\) &\\
\bottomrule
\end{tabular}
\end{table}

因为
\[
    \lambda_1=2
\]
为最大正检验数，所以 \(x_1\) 入基。按 \(\theta\) 规则，\(x_3\) 出基。换基运算为
\[
    R_{x_1}=R_{x_3},\qquad
    R_{x_4}\leftarrow R_{x_4}-2R_{x_1},\qquad
    R_{x_5}\leftarrow R_{x_5}-R_{x_1},\qquad
    R_{x_6}\leftarrow R_{x_6}-R_{x_1}.
\]

\begin{table}[H]
\centering
\zihao{-4}
\setlength{\tabcolsep}{2.2pt}
\caption{第二次换基后的单纯形表}
\begin{tabular}{ccccccccccc}
\toprule
\multicolumn{3}{c}{$c_j\rightarrow$} & \(2\) & \(3\) & \(0\) & \(0\) & \(0\) & \(0\) & \(0\) & \multirow{2}{*}{\(\theta_i\)}\\
\(C_B\) & 基 & \(b\) & \(x_1\) & \(x_2\) & \(x_3\) & \(x_4\) & \(x_5\) & \(x_6\) & \(x_7\) &\\
\midrule
\(2\) & \(x_1\) & \(1.5\)  & \(1\) & \(0\) & \(1\) & \(0\) & \(0\) & \(0\) & \(-6\) & --\\
\(0\) & \(x_4\) & \(4.25\) & \(0\) & \(0\) & \(-2\) & \(1\) & \(0\) & \(0\) & \(11\) & \(4.25/11\)\\
\(0\) & \(x_5\) & \(1.25\) & \(0\) & \(0\) & \(-1\) & \(0\) & \(1\) & \(0\) & \([5]\) & \(1.25/5=0.25\)\\
\(0\) & \(x_6\) & \(3.5\)  & \(0\) & \(0\) & \(-1\) & \(0\) & \(0\) & \(1\) & \(6\) & \(3.5/6\)\\
\(3\) & \(x_2\) & \(4.75\) & \(0\) & \(1\) & \(0\) & \(0\) & \(0\) & \(0\) & \(1\) & \(4.75/1=4.75\)\\
\multicolumn{3}{c}{检验数} & \(0\) & \(0\) & \(-2\) & \(0\) & \(0\) & \(0\) & \(9\) &\\
\bottomrule
\end{tabular}
\end{table}

因为
\[
    \lambda_7=9
\]
为最大正检验数，所以 \(x_7\) 入基。按 \(\theta\) 规则，\(x_5\) 出基。先将主元行除以 \(5\)，再消去 \(x_7\) 列的其余系数：
\[
    R_{x_7}=\frac{1}{5}R_{x_5},
\]
\[
    R_{x_1}\leftarrow R_{x_1}+6R_{x_7},\quad
    R_{x_4}\leftarrow R_{x_4}-11R_{x_7},\quad
    R_{x_6}\leftarrow R_{x_6}-6R_{x_7},\quad
    R_{x_2}\leftarrow R_{x_2}-R_{x_7}.
\]

\begin{table}[H]
\centering
\zihao{-4}
\setlength{\tabcolsep}{3pt}
\caption{第三次换基后的单纯形表}
\begin{tabular}{cccccccccc}
\toprule
\multicolumn{3}{c}{$c_j\rightarrow$} & \(2\) & \(3\) & \(0\) & \(0\) & \(0\) & \(0\) & \(0\)\\
\(C_B\) & 基 & \(b\) & \(x_1\) & \(x_2\) & \(x_3\) & \(x_4\) & \(x_5\) & \(x_6\) & \(x_7\)\\
\midrule
\(2\) & \(x_1\) & \(3\)    & \(1\) & \(0\) & \(-1/5\) & \(0\) & \(6/5\) & \(0\) & \(0\)\\
\(0\) & \(x_4\) & \(1.5\)  & \(0\) & \(0\) & \(1/5\)  & \(1\) & \(-11/5\) & \(0\) & \(0\)\\
\(0\) & \(x_7\) & \(0.25\) & \(0\) & \(0\) & \(-1/5\) & \(0\) & \(1/5\) & \(0\) & \(1\)\\
\(0\) & \(x_6\) & \(2\)    & \(0\) & \(0\) & \(1/5\)  & \(0\) & \(-6/5\) & \(1\) & \(0\)\\
\(3\) & \(x_2\) & \(4.5\)  & \(0\) & \(1\) & \(1/5\)  & \(0\) & \(-1/5\) & \(0\) & \(0\)\\
\multicolumn{3}{c}{检验数} & \(0\) & \(0\) & \(-1/5\) & \(0\) & \(-9/5\) & \(0\) & \(0\)\\
\bottomrule
\end{tabular}
\end{table}

表 1-4 中所有检验数均非正，所以当前基可行解为最优解：
\[
    X^*=(3,4.5,0,1.5,0,2,0.25)^{\mathrm T},\qquad z^*=19.5.
\]

\subsection{补充问题}

对于第一题(b)，合同要求
\[
    x_1\ge 3,\qquad x_2\ge 5. \tag{1-30}
\]
但水资源约束仍为
\[
    x_1+6x_2\le 30. \tag{1-31}
\]
由 \(x_2\ge 5\) 可知 \(6x_2\ge 30\)，再结合 \(x_1\ge 3\)，得
\[
    x_1+6x_2\ge 33>30.
\]
这与式(1-31)矛盾。因此合同要求不可行，农夫不应签订合同。

对于第一题(c)，若增加 \(10000\) gallons water，则水资源约束变为
\[
    x_1+6x_2\le 40. \tag{1-32}
\]
其余约束不变。由于 corn 的单位利润系数较大，先令
\[
    x_2=4.75.
\]
由资本约束 \(x_1+x_2\le 7.5\)，得
\[
    x_1\le 2.75.
\]
取
\[
    x_1=2.75.
\]
检查水与土地约束：
\[
    x_1+6x_2=2.75+6\cdot 4.75=31.25\le 40,
\]
\[
    2x_1+x_2=2\cdot 2.75+4.75=10.25\le 12.
\]
因此该方案可行，利润为
\[
    2\cdot 2.75+3\cdot 4.75=19.75,
\]
即 \(1975\) dollars。与原最优利润 \(1950\) dollars 相比，仅增加 \(25\) dollars，小于额外水资源成本 \(50\) dollars。因此，农夫不应购买额外水资源。

\section{第二题}

\subsection{问题叙述}

2. A factory plans to produce a seasonal product over the next 3 months. The monthly demand, production cost per unit, and storage cost per unit are given below:

\begin{center}
\begin{tabular}{ccc}
\toprule
Month \(t\) & Demand \(d_t\) & Production cost \(c_t\) (¥/unit)\\
\midrule
1 & 80  & 10\\
2 & 120 & 12\\
3 & 100 & 11\\
\bottomrule
\end{tabular}
\end{center}

\begin{itemize}[leftmargin=2.3em]
  \item Storage cost per unit per month: ¥1.
  \item Maximum production capacity per month: 150 units.
  \item Initial storage at the beginning of the \(1^{\mathrm{st}}\) month: 20 units.
  \item Required final storage at the end of the \(3^{\mathrm{rd}}\) month: no less than 10 units.
  \item Stockout is not allowed (all demand must be satisfied on time).
\end{itemize}

The factory wants to determine the monthly production quantities that minimize the total cost over the 3 months horizon. Please use simplex method.

\subsection{模型建立}

令
\[
    p_t=\text{第 }t\text{ 月的生产量},\qquad
    I_t=\text{第 }t\text{ 月末库存量}.
\]
已知初始库存为
\[
    I_0=20.
\]
库存平衡关系为
\[
    I_t=I_{t-1}+p_t-d_t,\qquad t=1,2,3. \tag{2-1}
\]
由于所有生产成本均为正，最优解中不会生产超过满足需求与最低期末库存所需的数量，因此有
\[
    I_3=10. \tag{2-2}
\]
由库存平衡关系可得
\[
    p_1=60+I_1, \tag{2-3}
\]
\[
    p_2=120+I_2-I_1, \tag{2-4}
\]
\[
    p_3=110-I_2. \tag{2-5}
\]
其中，\(I_1,I_2\ge 0\) 保证前两个月不缺货；式(2-2)保证第三个月末库存不低于 \(10\) units。

生产能力约束 \(p_t\le 150\) 给出
\[
    I_1\le 90, \tag{2-6}
\]
\[
    -I_1+I_2\le 30, \tag{2-7}
\]
\[
    I_2\le 110. \tag{2-8}
\]
其中 \(p_1,p_2,p_3\ge 0\) 在最终解中还需检查。

本文将保管费理解为跨入下一月的库存费用，因此总成本为
\[
    C=10p_1+12p_2+11p_3+I_1+I_2. \tag{2-9}
\]
若将第 \(3\) 月末的最低库存也计入一次期末保管费，则总成本只需在式(2-9)基础上再加 \(I_3=10\)，不会改变最优生产计划。
将式(2-3)、式(2-4)、式(2-5)代入式(2-9)，得
\[
\begin{aligned}
    C
    &=10(60+I_1)+12(120+I_2-I_1)+11(110-I_2)+I_1+I_2\\
    &=3250-I_1+2I_2.
\end{aligned} \tag{2-10}
\]
于是，最小化 \(C\) 等价于最大化
\[
    v=I_1-2I_2. \tag{2-11}
\]
第二题可化为标准的最大化模型：
\[
    \begin{aligned}
    \max\quad & v=I_1-2I_2,\\
    \mathrm{s.t.}\quad
      &I_1\le 90,\\
      &-I_1+I_2\le 30,\\
      &I_2\le 110,\\
      &I_1,I_2\ge 0.
    \end{aligned} \tag{2-12}
\]

\subsection{单纯形法推导}

第二题的变量较少，但换基逻辑与第一题完全一致：先找到一个显然可行的基解，再观察目标函数中非基变量的系数；正系数表示增加该变量能够改进目标函数，最小正比值则保证移动后仍然满足所有非负约束。

引入松弛变量 \(u_1,u_2,u_3\)，将式(2-12)化为
\[
    \begin{aligned}
    \max\quad & v=I_1-2I_2,\\
    \mathrm{s.t.}\quad
      &I_1+u_1=90,\\
      &-I_1+I_2+u_2=30,\\
      &I_2+u_3=110,\\
      &I_1,I_2,u_1,u_2,u_3\ge 0.
    \end{aligned} \tag{2-13}
\]
取 \(u_1,u_2,u_3\) 为初始基变量，\(I_1,I_2\) 为非基变量，则
\[
    \begin{cases}
    u_1=90-I_1,\\
    u_2=30+I_1-I_2,\\
    u_3=110-I_2.
    \end{cases} \tag{2-14}
\]
令非基变量 \(I_1=I_2=0\)，得初始基可行解
\[
    (I_1,I_2,u_1,u_2,u_3)=(0,0,90,30,110),\qquad v^{(0)}=0. \tag{2-15}
\]
目标函数为
\[
    v=I_1-2I_2. \tag{2-16}
\]
当前非基变量为 \(I_1,I_2\)。若把 \(I_1\) 从 \(0\) 增加 \(\Delta\)，则 \(v\) 增加 \(\Delta\)；若把 \(I_2\) 从 \(0\) 增加 \(\Delta\)，则 \(v\) 减少 \(2\Delta\)。因此只有 \(I_1\) 是能改善目标函数的方向，故选择 \(I_1\) 为入基变量。

保持 \(I_2=0\)，让 \(I_1\) 从 \(0\) 开始增加。由式(2-14)可得
\[
    u_1=90-I_1\ge 0,\qquad
    u_2=30+I_1\ge 0,\qquad
    u_3=110\ge 0.
\]
其中 \(u_2=30+I_1\) 会随 \(I_1\) 增加而增加，\(u_3\) 与 \(I_1\) 无关，因此它们都不会阻止 \(I_1\) 增大。唯一会随 \(I_1\) 增大而减小的是 \(u_1=90-I_1\)，其非负性要求
\[
    I_1\le 90.
\]
所以本次允许的最大步长为
\[
    \theta=\frac{90}{1}=90. \tag{2-17}
\]
当 \(I_1=90\) 时，\(u_1=0\)，故 \(u_1\) 为出基变量。

由
\[
    u_1=90-I_1
\]
解出
\[
    I_1=90-u_1. \tag{2-18}
\]
这一步把 \(I_1\) 改写为新的基变量，并把 \(u_1\) 改写为新的非基变量。随后令新的非基变量 \(u_1=I_2=0\)，即可得到换基后的顶点。

代入式(2-14)，得
\[
    \begin{cases}
    I_1=90-u_1,\\
    u_2=120-u_1-I_2,\\
    u_3=110-I_2.
    \end{cases} \tag{2-19}
\]
目标函数变为
\[
    v=90-u_1-2I_2. \tag{2-20}
\]
式(2-20)中，当前非基变量 \(u_1,I_2\) 的系数分别为 \(-1\) 和 \(-2\)。由于非基变量只能从 \(0\) 向正方向增加，任何增加都会使 \(v\) 减小，因而不存在继续改进目标函数的相邻方向。因此当前基可行解已为最优解：
\[
    I_1^*=90,\qquad I_2^*=0,\qquad v^*=90. \tag{2-21}
\]
由式(2-3)、式(2-4)、式(2-5)，生产计划为
\[
    p_1=60+90=150,
\]
\[
    p_2=120+0-90=30,
\]
\[
    p_3=110-0=110.
\]
总成本为
\[
    C^*=10\cdot150+12\cdot30+11\cdot110+90+0=3160.
\]
若将第 \(3\) 月末库存也计入一次期末保管费，则总成本为 \(3170\)，生产计划不变。
因此第二题的最优生产计划为
\[
    \boxed{p_1=150,\qquad p_2=30,\qquad p_3=110,}
\]
最小总成本为
\[
    \boxed{3160\text{ yuan}.}
\]

\subsection{单纯形表}

第二题的单纯形表同样按“最大正检验数确定入基列、最小正比值确定出基行”的规则计算。由于只有 \(I_1\) 的检验数为正，入基变量唯一；再由入基列中正系数对应的比值确定出基变量。

\begin{table}[H]
\centering
\zihao{-4}
\setlength{\tabcolsep}{4pt}
\caption{初始单纯形表}
\begin{tabular}{ccccccccc}
\toprule
\multicolumn{3}{c}{$c_j\rightarrow$} & \(1\) & \(-2\) & \(0\) & \(0\) & \(0\) & \multirow{2}{*}{\(\theta_i\)}\\
\(C_B\) & 基 & \(b\) & \(I_1\) & \(I_2\) & \(u_1\) & \(u_2\) & \(u_3\) &\\
\midrule
\(0\) & \(u_1\) & \(90\)  & \([1]\) & \(0\) & \(1\) & \(0\) & \(0\) & \(90/1=90\)\\
\(0\) & \(u_2\) & \(30\)  & \(-1\) & \(1\) & \(0\) & \(1\) & \(0\) & --\\
\(0\) & \(u_3\) & \(110\) & \(0\) & \(1\) & \(0\) & \(0\) & \(1\) & --\\
\multicolumn{3}{c}{检验数} & \(1\) & \(-2\) & \(0\) & \(0\) & \(0\) &\\
\bottomrule
\end{tabular}
\end{table}

表 2-1 中，\(\lambda_{I_1}=1\) 为最大正检验数，故 \(I_1\) 入基；只有第一行对应正系数，故 \(u_1\) 出基。换基运算为
\[
    R_{I_1}=R_{u_1},\qquad R_{u_2}\leftarrow R_{u_2}+R_{I_1}.
\]

\begin{table}[H]
\centering
\zihao{-4}
\setlength{\tabcolsep}{4pt}
\caption{换基后的单纯形表}
\begin{tabular}{cccccccc}
\toprule
\multicolumn{3}{c}{$c_j\rightarrow$} & \(1\) & \(-2\) & \(0\) & \(0\) & \(0\)\\
\(C_B\) & 基 & \(b\) & \(I_1\) & \(I_2\) & \(u_1\) & \(u_2\) & \(u_3\)\\
\midrule
\(1\) & \(I_1\) & \(90\)  & \(1\) & \(0\) & \(1\) & \(0\) & \(0\)\\
\(0\) & \(u_2\) & \(120\) & \(0\) & \(1\) & \(1\) & \(1\) & \(0\)\\
\(0\) & \(u_3\) & \(110\) & \(0\) & \(1\) & \(0\) & \(0\) & \(1\)\\
\multicolumn{3}{c}{检验数} & \(0\) & \(-2\) & \(-1\) & \(0\) & \(0\)\\
\bottomrule
\end{tabular}
\end{table}

表 2-2 中所有检验数均非正，因此得到最优解
\[
    I_1^*=90,\qquad I_2^*=0.
\]
相应生产计划为
\[
    p_1=150,\qquad p_2=30,\qquad p_3=110.
\]
检查库存：
\[
    I_1=20+150-80=90,
\]
\[
    I_2=90+30-120=0,
\]
\[
    I_3=0+110-100=10.
\]
所有需求均按时满足，且第三个月末库存达到 \(10\) units。

\section*{总结}

第一题(a)的最优方案为种植 \(300\) bushels tomatoes 与 \(450\) bushels corn，最大利润为 \(1950\) dollars。第一题(b)的合同要求与水资源约束矛盾，因此不应签订。第一题(c)购买额外水资源只能增加 \(25\) dollars 的利润，但需支付 \(50\) dollars 的成本，因此不应购买。

第二题的最优生产计划为
\[
    (p_1,p_2,p_3)=(150,30,110),
\]
相应库存为
\[
    (I_1,I_2,I_3)=(90,0,10),
\]
最小总成本为 \(3160\) yuan。
若将第 \(3\) 月末库存也计入一次期末保管费，则总成本为 \(3170\) yuan，最优生产计划不变。

\end{document}

